\src{188}
\begin{equation}\tag{7.2}
  f^{*}:H^{*}(X)\longrightarrow H^{*}(Y).
\end{equation}
Since these are Poincaré algebras, transposition gives a homomorphism of vector
spaces
\begin{equation}\tag{7.3}
  f_{*}:H^{*}(Y)\longrightarrow H^{*}(X).
\end{equation}
We saw in Section 4 that the effect of $f^{*}$ on the cohomology classes
corresponding to nonsingular cycles has the geometric interpretation of taking
the cohomology classes corresponding to their inverse images. In the present
setting it is important to show that (7.3) similarly corresponds to the
``direct image'' operation on cycles. Under suitable nonsingularity conditions,
at least, this follows from the following special case.

\result{THEOREM}{4} If $f$ is the inclusion of a nonsingular subvariety $Y^{m}$
of $X^{n}$ into $X^{n}$, then, writing $1_Y$ for the unit element of $H(Y)$,
one has
\begin{equation}\tag{7.4}
  f_{*}(1_Y)=P_X(Y),
\end{equation}
where the right-hand side is the cohomology class in $X$ associated with $Y$.

This formula is equivalent to
\begin{equation}\tag{7.4 bis}
  \bigl\langle\xi^{m,m},P_X(Y)\bigr\rangle\varepsilon_Y
  =f^{*}(\xi^{m,m})
  \qquad
  \bigl(\xi^{m,m}\in H^{m}(X,\Omega_X^{m})\bigr),
\end{equation}
\footnote{The printed text has $f_{*}(\xi^{m,m})$. Since
$\xi^{m,m}\in H^{m}(X,\Omega_X^{m})$ and both sides belong to
$H^{m}(Y,\Omega_Y^{m})$, the required map is $f^{*}$.} where $\varepsilon_Y$ is
the fundamental element in $H^{m}(Y,\Omega_Y^{m})$. For nonsingular projective
varieties this gives a new definition of the cohomology class associated with
$Y$. To prove Theorem 4, Theorem 3 lets us interpret the transpose of the
homomorphism
\[
  H^{m}(X,\Omega_X^{m})
  \longrightarrow H^{m}(Y,\Omega_Y^{m})
  =H^{m}(X,\Omega_Y^{m})
\]
as the homomorphism
\begin{equation}\tag{7.5}
\begin{split}
  H^{n-m}(X,\Omega_X^{n-m})
  &\simeq
  \operatorname{Ext}_{\Osh_X}^{n-m}(X;\Omega_X^{m},\Omega_X^{n})\\
  &\longleftarrow
  \operatorname{Ext}_{\Osh_X}^{n-m}(X;\Omega_Y^{m},\Omega_X^{n})
  \simeq
  \operatorname{Hom}_{\Osh_X}(X;\Omega_Y^{m},\Omega_Y^{m}).
\end{split}
\end{equation}
One checks that the element $1_Y$ of the dual of
$H^{m}(Y,\Omega_Y^{m})$ identifies with the element of the right-hand side
corresponding to the identity endomorphism of $\Omega_Y^{m}$, and also that the
image of this element in $H^{n-m}(X,\Omega_X^{n-m})$ is indeed $P_X(Y)$. These
compatibilities, which could have been stated already in Section 4, can be
formulated---and in fact hold---for arbitrary nonsingular varieties, the second

